%! Constructing a Confidence Interval for a Population Proportion — Unit 6, Lesson 6.2 — Exit Ticket
\documentclass[10pt]{article}
\usepackage{apstats-article}
\usepackage{apstats-boxes}
\newcommand{\TallMath}[1]{$\displaystyle #1\rule[-1.4em]{0pt}{3.2em}$}

\begin{document}

\pageheader{Unit 6, Lesson 6.2}{Exit Ticket}
\namedateperiod

\vspace{0.15in}

A random sample of 180 adult dog owners found that 126 say their dog sleeps in the bedroom.
Construct a 95\% confidence interval for the true proportion of adult dog owners whose dog
sleeps in the bedroom.

\begin{enumerate}
  \item \textbf{State} the parameter of interest.
        \writelines{1}

  \item \textbf{Plan}: Verify all three conditions.
        \begin{itemize}
          \item Random: \blank{7cm}
          \item Large Counts: $n\hat p = $ \blank{2.5cm} \quad $n(1-\hat p) = $ \blank{2.5cm} \quad Both $\ge 10$? \blank{1cm}
          \item 10\% rule: \blank{6cm}
        \end{itemize}

  \item \textbf{Do}: Compute the interval. ($\hat p = $ \blank{2cm})

        \[
          \hat p \pm z^* \sqrt{\frac{\hat p(1-\hat p)}{n}}
          = \blank{1.5cm} \pm \blank{1.5cm}\sqrt{\frac{\blank{2cm}}{\blank{1.5cm}}}
          = (\blank{2.5cm},\; \blank{2.5cm})
        \]

  \item \textbf{Conclude}: Write a complete interpretation.
        \writelines{2}
\end{enumerate}

\end{document}
